% year: תש"פ
% type: מבחן
% semester: א
% moed: ב
% part: ב
% number: 5
% points: 9
% categories: func-analysis, derivatives, multiple-choice
% source: exams/מבחנים/תשפ/מועד ב תשפ.pdf

\begin{question}
אם $f(x)$ פונקציה ממשית שנגזרתה היא $f'(x)=(x-1)^2(x+2)(x-5)$ אזי:
\begin{enumerate}
  \item ל-$f(x)$ יש מקסימום ב-$x=-2$ ומינימום ב-$x=5$, ואין נק' קיצון ב-$x=1$.
  \item ל-$f(x)$ יש מקסימום ב-$x=5$ ומינימום ב-$x=-2$, ואין נק' קיצון ב-$x=1$.
  \item ל-$f(x)$ יש מקסימום ב-$x=-2$ וב-$x=1$ ומינימום ב-$x=5$.
  \item ל-$f(x)$ יש מקסימום ב-$x=5$ ומינימום ב-$x=-2$ וב-$x=1$.
  \item ל-$f(x)$ אין נקודות קיצון.
\end{enumerate}
\end{question}

\begin{hint}
בדקו את סימן הנגזרת משני צידי כל נקודה קריטית. שימו לב שהגורם $(x-1)^2$ אינו מחליף סימן.
\end{hint}

\begin{solution}
\textbf{התשובה הנכונה: א.}

$f$ גזירה בכל $\R$, ולכן נקודות קיצון אפשריות רק היכן ש-$f'=0$ (משפט פרמה): $x=-2,\,1,\,5$. נבדוק את סימן $f'$. הגורם $(x-1)^2\ge0$ ואינו מחליף סימן, ולכן הסימן נקבע על ידי $(x+2)(x-5)$:
\begin{itemize}
  \item $x<-2$: $(x+2)(x-5)>0$, ולכן $f'>0$ --- $f$ עולה.
  \item $-2<x<5$, $x\neq 1$: $(x+2)(x-5)<0$, ולכן $f'<0$ --- $f$ יורדת.
  \item $x>5$: $f'>0$ --- $f$ עולה.
\end{itemize}
לכן ב-$x=-2$ הנגזרת עוברת מחיובית לשלילית --- מקסימום מקומי; ב-$x=5$ היא עוברת משלילית לחיובית --- מינימום מקומי; וב-$x=1$ הנגזרת שלילית משני הצדדים, כך ש-$f$ יורדת ממש בסביבת $1$ ואין שם קיצון.
\end{solution}
